%
% listener_lifecycle.tex
%
% A Z specification of the Hub display link's listener/connection lifecycle:
% one reconnecting connection, a reader thread per dial that stop() may fail to
% join, the reply sinks those threads feed, and the report a reader makes when
% its stream breaks.  This is the FIXED design: a connection generation owns
% its socket, its reply sink and the right to be closed by its own reader;
% a reader that outlives its generation can touch nothing of the next.
%
% Models the defect reported twice on the lux-ghhf branch (round 2 L-2..L-4;
% round 3 R3-1 and SF3-2) and the design that closes it.  The code it governs:
% src/punt_lux/domain/hub/{listener_thread,display_replies,display_connection,
% inbound_router,display_link}.py.
%
% Type-check:  fuzz docs/listener_lifecycle.tex
% Model-check (see the Verification section for every goal predicate):
%   probcli docs/listener_lifecycle.tex -model_check -p DEFAULT_SETSIZE 2 \
%     -p MAX_OPERATIONS 200
%
% Fidelity controls (each differs from this file in the stated place only):
%   docs/listener_lifecycle_current_buggy.tex    the code as it stands
%   docs/listener_lifecycle_gate_only_buggy.tex  the stopping flag as the only fix
%

\documentclass[a4paper,10pt,fleqn]{article}
\usepackage[margin=1in]{geometry}
\usepackage{fuzz}
\usepackage[colorlinks=true,linkcolor=blue,citecolor=blue,urlcolor=blue]{hyperref}

\begin{document}

\title{Listener Lifecycle: a Z Specification}
\author{Formal model of the display link's reader, its socket, and its replies}
\date{October 2026}
\maketitle

% ============================================================
\section{Introduction}
% ============================================================

\texttt{DisplayConnection} is one object that holds, at any moment, at most
one socket to the display and is re-opened after every loss.  A
\texttt{ListenerThread} reads one socket on a daemon thread and hands what it
decodes to the \texttt{InboundRouter}; when its stream breaks it calls
\texttt{DisplayReplies.end}, which closes \emph{whatever the connection holds
now} and empties the reply queues.  \texttt{ListenerThread.stop} joins with a
two-second timeout, and \texttt{start} deliberately abandons a thread it could
not join.  Two review rounds in succession found the consequence, from two
sides:

\begin{itemize}
  \item \textbf{R3-1.}\ An abandoned reader whose read fails after the
    reconnect finds the connection open --- the \emph{replacement} is open ---
    and reports a broken stream.  The healthy replacement is closed, and the
    replies queued for it are dropped.
  \item \textbf{SF3-2.}\ An abandoned reader goes on routing the rest of its
    old buffer into the shared queues.  A stale \texttt{Ack},
    \texttt{Pong} or \texttt{QueryResponse} fills the request of a caller that
    asked on the new socket, and a torn frame in that buffer reports the wrong
    connection broken.
\end{itemize}

Round~2 found the same class (L-2..L-4).  Two rounds of the same class is the
recurrence signal: the problem is a state space, and this file checks it.
Both findings are one defect.  \emph{A decision about generation~$g$ is
executed against generation $g' \ne g$} because nothing in the call carries
the generation: \texttt{is\_open} describes the current socket, and
\texttt{end} acts on the current socket.  The reviewers' suggested repair ---
consult the thread's \texttt{stopping} event before acting --- narrows the
window and leaves it open (Section~\ref{sec:fidelity}); it is checked here as
the second fidelity control, not assumed.

\paragraph{The design, by responsibility.}
The socket's lifetime is owned by exactly one thing: the dial that made it.
The fixed design says so in the type.  Every dial produces a \emph{generation}
(in the code, a per-socket object: socket, endpoint, ready message, reply
sink, and the reader started on it).  Three consequences, each a clause of
the model:

\begin{description}
  \item[Routing is keyed by generation.] A reader holds the sink of its own
    generation and nothing else.  A reply from generation~$g$ goes to the sink
    of~$g$; when $g$ is closed that sink is discarded, and a late put into it
    is dropped.  There is no check-then-put to tear, because the sink the
    reader holds \emph{is} the check.  (\texttt{RouteLive}, \texttt{RouteDead}.)
  \item[A report is a compare-and-close on the reader's own generation.]
    ``My stream broke'' closes generation $g$ if $g$ is still the connection's
    current generation and open, and is a no-op otherwise.  (\texttt{BrokenOwn},
    \texttt{BrokenStale}.)  This is the only operation a reader has on the
    connection, and it takes the reader's generation as an argument.
  \item[The \texttt{stopping} flag silences, it does not protect.] Once set,
    the reader takes no further step that reports or routes, so a stopped
    reader neither logs a spurious warning nor does stray work.  It is a
    courtesy; the two invariants above hold without it
    (Section~\ref{sec:fidelity}, second control).
\end{description}

A thread that cannot be joined is therefore abandoned \emph{together with its
socket}: the generation is retired (closed, sink discarded) before a
successor is started on the connection.  A successor never reads a socket an
abandoned thread is still reading (L4).

\paragraph{Decisions this model rests on; each is the implementer's duty.}
\begin{enumerate}
  \item \emph{The compare-and-close and the dial's publication are atomic with
    respect to each other.}  The code needs a small state lock, held only for
    the field update and never across I/O or a dial.  It must \emph{not} be
    the connection's \texttt{\_lock}: that lock is held across
    \texttt{sendall} and \texttt{recv}, and a sender stuck on a wedged display
    would then hold up the report that would have closed the wedged socket
    (the very stall that makes \texttt{stop} time out in R3-1).  See
    ``Lock order'' below.
  \item \emph{A reply sink belongs to a generation.}  The ack, pong and
    query-response queues are per generation.  The business-event queue and
    the interaction handler are \emph{link}-level and are outside L1: an event
    the display emitted is real whichever socket it left on, and it
    is not rerouted.
  \item \emph{Abandonment retires the socket.}  Before a successor reader is
    started over a reader that did not join, the generation is closed, by
    \texttt{shutdown} then \texttt{close}, so that the abandoned thread's
    blocked read returns.
\end{enumerate}

\paragraph{Scope.}
The carrier is three generations (two redials past the first dial), two
reader threads (an abandoned reader and its successor; the code never has
more live threads than that, because it abandons one at a time), one buffered
frame per reader, and a reply as a bare tag naming its generation.  Frame
contents, the handler's duration and the byte-level decode are abstracted:
a read failure, an end-of-file and a malformed frame are one event, because
the code answers all three by the same call (\texttt{end}), and
\texttt{end} is a no-op on a closed connection.  The thread running a handler
for longer than the join timeout is exactly a reader with a non-empty buffer
that \texttt{StopTimeout} abandons.

\paragraph{Invariants.}

\begin{description}
  \item[L1 --- replies are routed within their generation.] No reply is
    queued except in the sink of the live current generation; no waiter
    exists on a generation that is not live.  (A stale tag in the queue is
    what a fresh waiter consumes as the answer to a question it never asked.)
  \item[L2 --- a report never closes another generation.] A generation is
    closed only deliberately (the link closing it) or because its own reader
    found its own stream broken.
  \item[L3 --- a stopped reader is silent.] No routing or reporting step is
    taken by a reader whose stopping flag is set and which holds no
    permission granted before it was set.
  \item[L4 --- one reader per live generation.] No generation is read by two
    threads at once.
\end{description}

% ============================================================
\section{Constants and Free Types}
% ============================================================

\begin{axdef}
  maxGen : \nat \\
  maxBuf : \nat
\where
  maxGen = 3 \\
  maxBuf = 1
\end{axdef}

Three generations: the original, the replacement, and one more, so that a
replacement can itself be replaced while an abandoned reader of the first is
still alive.  One buffered frame is the least that lets an abandoned reader
have something left to route.

\begin{zed}
  T ::= ta | tb
\end{zed}

The two reader threads.  A thread is a slot: it is idle (\texttt{lexit}),
running, or stopping.

\begin{zed}
  LST ::= lexit | lrun | lstop
\end{zed}

\texttt{lexit}: no thread.  \texttt{lrun}: the thread runs and its stopping
flag is clear.  \texttt{lstop}: the flag is set and the thread is alive ---
\texttt{stop} timed out, or has not yet joined.  This is the abandoned thread.

\begin{zed}
  FSTATE ::= fclear | fset
\end{zed}

% ============================================================
\section{State}
% ============================================================

The state is split into three groups, then joined, so that an operation can
say what it leaves alone.

\begin{schema}{Conn}
  curGen : \nat \\
  held : FSTATE \\
  waiters : \power \nat \\
  qgen : \power \nat
\where
  curGen \leq maxGen \\
  waiters \subseteq 0 \upto maxGen \\
  qgen \subseteq 0 \upto maxGen \\
  held = fset \implies curGen > 0
\end{schema}

\texttt{curGen} is the latest generation dialed (zero before the first dial);
\texttt{held} says whether the connection holds that generation's socket ---
\texttt{is\_open}.  Every earlier generation is closed.  \texttt{waiters} are
the generations on which a caller is blocked awaiting a reply.
\texttt{qgen} are the generations that have a reply sitting in a sink (a bare
tag per reply; one suffices, since duplicates change nothing below).

\begin{schema}{Thr}
  lst : T \fun LST \\
  lgen : T \fun \nat \\
  armed : T \fun FSTATE \\
  buf : T \fun \nat
\where
  \forall t : T \spot lgen~t \leq maxGen \land buf~t \leq maxBuf
\end{schema}

\texttt{lgen} is the generation a thread was started on (meaningful while the
thread is not \texttt{lexit}); \texttt{buf} counts frames it has read and not
yet routed.  \texttt{armed} is the permission the thread obtained the last time
it looked at its stopping flag while the flag was clear: the loop-top test in
\texttt{\_run}, and the re-test this design adds before a route or a report.
A step that routes or reports consumes the permission; the next needs a
fresh look.  The flag can be set \emph{between} the look and the step, which
is why a permission is state: it is the window, and the model keeps it open.

\begin{schema}{Hist}
  broken : \power \nat \\
  deliberate : \power \nat \\
  leak : FSTATE
\where
  broken \subseteq 0 \upto maxGen \\
  deliberate \subseteq 0 \upto maxGen
\end{schema}

Three history variables, never read by an operation of this specification, so
that the invariants are statements about what happened.  \texttt{broken}: the
generations whose own reader found their own stream broken.
\texttt{deliberate}: the generations the link closed.  \texttt{leak}: set when
a step is taken by a stopped reader holding no permission.

\begin{schema}{LL}
  Conn \\
  Thr \\
  Hist
\end{schema}

\section{Initialization}

\begin{schema}{Init}
  LL'
\where
  curGen' = 0 \\
  held' = fclear \\
  waiters' = \emptyset \\
  qgen' = \emptyset \\
  lst' = \{ ta \mapsto lexit, tb \mapsto lexit \} \\
  lgen' = \{ ta \mapsto 0, tb \mapsto 0 \} \\
  armed' = \{ ta \mapsto fclear, tb \mapsto fclear \} \\
  buf' = \{ ta \mapsto 0, tb \mapsto 0 \} \\
  broken' = \emptyset \\
  deliberate' = \emptyset \\
  leak' = fclear
\end{schema}

Nothing dialed, no reader, no reply, nothing recorded.

% ============================================================
\section{Operations on the Connection}
% ============================================================

\subsection{Dial}

\begin{schema}{Dial}
  \Delta LL \\
  \Xi Thr \\
  \Xi Hist
\where
  held = fclear \\
  curGen < maxGen \\
  curGen' = curGen + 1 \\
  held' = fset \\
  waiters' = waiters \\
  qgen' = qgen
\end{schema}

\texttt{DisplayConnection.open}: a new generation, a new socket.  Only a
closed connection dials.  The new generation has no sink content and no
waiter; nothing of an earlier generation is carried over.

\subsection{Close}

\begin{schema}{Close}
  \Delta LL \\
  \Xi Thr
\where
  held = fset \\
  held' = fclear \\
  curGen' = curGen \\
  waiters' = \emptyset \\
  qgen' = \emptyset \\
  deliberate' = deliberate \cup \{ curGen \} \\
  broken' = broken \\
  leak' = leak
\end{schema}

The link closes the connection on purpose: \texttt{DisplayLink.close},
\texttt{ClientRegistry.drop}.  The generation's sink is discarded and its
waiters fail.  Readers are not touched; \texttt{StopJoined} and
\texttt{StopTimeout} are separate steps, in either order.

\subsection{Request and Consume}

\begin{schema}{Request}
  \Delta LL \\
  \Xi Thr \\
  \Xi Hist
\where
  held = fset \\
  curGen \notin waiters \\
  waiters' = waiters \cup \{ curGen \} \\
  curGen' = curGen \\
  held' = held \\
  qgen' = qgen
\end{schema}

A caller sends a ping, a query or a scene on the current socket and blocks for
the reply.

\begin{schema}{Consume}
  \Delta LL \\
  \Xi Thr \\
  \Xi Hist
\where
  curGen \in waiters \cap qgen \\
  waiters' = waiters \setminus \{ curGen \} \\
  qgen' = qgen \setminus \{ curGen \} \\
  curGen' = curGen \\
  held' = held
\end{schema}

The blocked caller takes the reply from the sink of its own generation.

% ============================================================
\section{Operations on a Reader}
% ============================================================

Each takes the thread \texttt{t?}.

\subsection{StartListener}

\begin{schema}{StartListener}
  \Delta LL \\
  \Xi Conn \\
  \Xi Hist \\
  t? : T
\where
  lst~t? = lexit \\
  held = fset \\
  \forall u : T \spot lst~u \neq lrun \\
  \forall u : T \spot lst~u = lexit \lor lgen~u \neq curGen \\
  lst' = lst \oplus \{ t? \mapsto lrun \} \\
  lgen' = lgen \oplus \{ t? \mapsto curGen \} \\
  armed' = armed \oplus \{ t? \mapsto fclear \} \\
  buf' = buf \oplus \{ t? \mapsto 0 \}
\end{schema}

\texttt{ListenerThread.start}.  It is a no-op while a running reader exists
(third line); that is the code's \texttt{if not stopping: return}.  The
fourth line is this design's addition: a successor is started only on a
generation no live thread is still reading.  Abandoning a reader therefore
means retiring its generation first (Decision~3).

\subsection{CheckPermit}

\begin{schema}{CheckPermit}
  \Delta LL \\
  \Xi Conn \\
  \Xi Hist \\
  t? : T
\where
  lst~t? = lrun \\
  armed~t? = fclear \\
  lst' = lst \\
  lgen' = lgen \\
  buf' = buf \\
  armed' = armed \oplus \{ t? \mapsto fset \}
\end{schema}

The reader looks at its stopping flag, finds it clear, and obtains permission
for its next routing or reporting step.  Only a running reader can: once the
flag is set, no permission is issued.

\subsection{Feed}

\begin{schema}{Feed}
  \Delta LL \\
  \Xi Conn \\
  \Xi Hist \\
  t? : T
\where
  lst~t? = lrun \\
  held = fset \\
  lgen~t? = curGen \\
  buf~t? < maxBuf \\
  lst' = lst \\
  lgen' = lgen \\
  armed' = armed \\
  buf' = buf \oplus \{ t? \mapsto buf~t? + 1 \}
\end{schema}

The reader reads a chunk off its own, open socket.  A reader over a closed or
replaced socket reads nothing new; what it already holds stays in \texttt{buf}.

\subsection{Stop, timeout, and exit}

\begin{schema}{StopJoined}
  \Delta LL \\
  \Xi Conn \\
  \Xi Hist \\
  t? : T
\where
  lst~t? = lrun \\
  lst' = lst \oplus \{ t? \mapsto lexit \} \\
  lgen' = lgen \\
  armed' = armed \oplus \{ t? \mapsto fclear \} \\
  buf' = buf
\end{schema}

\texttt{stop}: the flag is set and the join succeeds.

\begin{schema}{StopTimeout}
  \Delta LL \\
  \Xi Conn \\
  \Xi Hist \\
  t? : T
\where
  lst~t? = lrun \\
  lst' = lst \oplus \{ t? \mapsto lstop \} \\
  lgen' = lgen \\
  armed' = armed \\
  buf' = buf
\end{schema}

\texttt{stop}: the flag is set and the join times out.  The thread is alive,
its permission (if any) intact, and it is abandoned.

\begin{schema}{ExitStopped}
  \Delta LL \\
  \Xi Conn \\
  \Xi Hist \\
  t? : T
\where
  lst~t? = lstop \\
  lst' = lst \oplus \{ t? \mapsto lexit \} \\
  lgen' = lgen \\
  armed' = armed \oplus \{ t? \mapsto fclear \} \\
  buf' = buf
\end{schema}

The abandoned thread sees its flag at its next test and exits.

\subsection{Route}

\begin{schema}{RouteLive}
  \Delta LL \\
  \Xi Hist \\
  t? : T
\where
  armed~t? = fset \\
  buf~t? > 0 \\
  held = fset \\
  lgen~t? = curGen \\
  lst' = lst \\
  lgen' = lgen \\
  armed' = armed \oplus \{ t? \mapsto fclear \} \\
  buf' = buf \oplus \{ t? \mapsto buf~t? - 1 \} \\
  curGen' = curGen \\
  held' = held \\
  waiters' = waiters \\
  qgen' = qgen \cup \{ lgen~t? \}
\end{schema}

\begin{schema}{RouteDead}
  \Delta LL \\
  \Xi Conn \\
  \Xi Hist \\
  t? : T
\where
  armed~t? = fset \\
  buf~t? > 0 \\
  \lnot (held = fset \land lgen~t? = curGen) \\
  lst' = lst \\
  lgen' = lgen \\
  armed' = armed \oplus \{ t? \mapsto fclear \} \\
  buf' = buf \oplus \{ t? \mapsto buf~t? - 1 \}
\end{schema}

A reader with permission routes one buffered reply into the sink of its own
generation.  If that generation is the live current one, the reply is queued
(\texttt{RouteLive}); if the generation has been closed or replaced, its sink
is gone and the reply is dropped (\texttt{RouteDead}).  Nothing in
\texttt{RouteDead} reaches the current generation.

\subsection{Broken}

\begin{schema}{BrokenOwn}
  \Delta LL \\
  t? : T
\where
  armed~t? = fset \\
  held = fset \\
  lgen~t? = curGen \\
  held' = fclear \\
  curGen' = curGen \\
  waiters' = \emptyset \\
  qgen' = \emptyset \\
  broken' = broken \cup \{ curGen \} \\
  deliberate' = deliberate \\
  leak' = leak \\
  lst' = lst \oplus \{ t? \mapsto lexit \} \\
  armed' = armed \oplus \{ t? \mapsto fclear \} \\
  lgen' = lgen \\
  buf' = buf
\end{schema}

\begin{schema}{BrokenStale}
  \Delta LL \\
  \Xi Conn \\
  \Xi Hist \\
  t? : T
\where
  armed~t? = fset \\
  \lnot (held = fset \land lgen~t? = curGen) \\
  lst' = lst \oplus \{ t? \mapsto lexit \} \\
  armed' = armed \oplus \{ t? \mapsto fclear \} \\
  lgen' = lgen \\
  buf' = buf
\end{schema}

The reader's stream broke --- the read failed, the display closed it, or a
frame did not decode --- and it reports.  The report is a compare-and-close
on the reader's own generation.  If that generation is the open current one,
it is closed and recorded as broken (\texttt{BrokenOwn}); otherwise the report
is a no-op (\texttt{BrokenStale}) and the reader exits.  The current
generation is never touched on behalf of another's reader.

\subsection{Halt}

\begin{schema}{Halt}
  \Xi LL
\where
  curGen = maxGen \\
  held = fclear \\
  \forall u : T \spot lst~u = lexit
\end{schema}

The carrier is finite, so the run ends when every generation has been dialed
and closed and every reader has exited.  \texttt{Halt} is enabled in that
quiescent state and nowhere else: a reachable state with no enabled operation
that is \emph{not} this one is a deadlock.

% ============================================================
\section{Invariants}
\label{sec:inv}
% ============================================================

Each is checked as the unreachability of its negation
(Section~\ref{sec:verify}); the negations are stated here as predicates over
the state.

\paragraph{L1 --- routing stays within a generation.}
\[ qgen \subseteq live \land waiters \subseteq live,
   \quad live = \{curGen\}\ \text{if}\ held = fset,\ \emptyset\ \text{otherwise.} \]
\texttt{RouteLive} adds only the reader's own generation and only when that
generation is the live current one; \texttt{RouteDead} adds nothing;
\texttt{Close} and \texttt{BrokenOwn} empty both sets; \texttt{Dial} and
\texttt{Request} add the new current generation only.  Hence both sets stay
within $live$.  A waiter on the new generation can therefore be answered only
from the new generation's sink.

\paragraph{L2 --- a generation is closed only by the link or by itself.}
\[ held = fclear \land curGen > 0 \implies curGen \in broken \cup deliberate. \]
\texttt{Close} records the generation it closes in \texttt{deliberate};
\texttt{BrokenOwn} records the generation it closes in \texttt{broken}, and
fires only when the reader's generation \emph{is} the current one.
\texttt{BrokenStale}, the only other report, leaves \texttt{Conn} unchanged.
No third operation clears \texttt{held}.

\paragraph{L3 --- a stopped reader is silent.}
\[ leak = fclear. \]
\texttt{leak} is set by a step that a stopped reader takes with no
permission.  Every routing and reporting step of this specification requires
permission, and \texttt{CheckPermit} issues one only to a \texttt{lrun}
reader.  A reader stopped \emph{after} it obtained permission may take that
one step --- the flag was clear when it looked, and the step is confined by
L1 and L2 whatever the flag says --- and no more.  This is a property of the
guards and is, deliberately, the weakest of the four: the first fidelity
control shows what its absence costs, and the second what it cannot buy.

\paragraph{L4 --- one reader per generation.}
\[ lst~ta \ne lexit \land lst~tb \ne lexit \land lgen~ta = lgen~tb. \]
\texttt{StartListener} refuses a generation that a live thread is still
reading.  (The negation is stated for two live threads on one generation; with
two slots that is the whole property.)

\paragraph{Deadlock-freedom, and the lock order the model assumes.}
The model has no lock-acquisition steps, because it makes the two critical
sections it needs atomic: the dial's publication of $(curGen, held)$ and the
compare-and-close both execute under one state lock, and a put into a sink
is a single queue operation.  What the model can and does check is that
\emph{no state other than the quiescent one lacks an enabled step}: every
reader can always exit or be stopped, every connection can always be closed
or dialed until the carrier is spent.  The lock discipline the
implementation owes the model is a graph with no nesting at all:

\begin{center}
\begin{tabular}{lll}
  lock & guards & held across \\ \hline
  state lock (new) & $curGen$, $held$, sink discard & field updates only; never I/O, never a dial \\
  \texttt{\_lock} (existing) & the socket's send and recv & one \texttt{sendall} or one \texttt{recv} \\
  reply-queue locks & one queue's contents & one \texttt{put} or \texttt{get} \\
\end{tabular}
\end{center}

No path takes one of these while holding another: the report takes the state
lock alone; the sender takes \texttt{\_lock} alone; the router takes a queue
lock alone.  A graph without edges has no cycle.  The one thing that would
create an edge is the tempting implementation of the compare-and-close under
\texttt{\_lock}: it would put the state update behind a wedged
\texttt{sendall}, and then the report waits on the sender, the sender waits on
a display that the report was about to abandon.  Decision~1 forbids it.

% ============================================================
\section{Fidelity Controls}
\label{sec:fidelity}
% ============================================================

\paragraph{The code as it stands ---
\texttt{listener\_lifecycle\_current\_buggy.tex}.}
The reply sink is one queue shared across reconnects (a reply is queued
whatever its generation); a report is \texttt{end}, which closes the current
connection if it is open, whatever generation the reporter read; a reader's
steps need only that the thread is alive, because \texttt{\_routed} tests the
flag at the top of the loop and not again; and \texttt{start} abandons a live
thread without retiring its socket.  ProB returns, for L2, an R3-1 trace in
ten steps: dial; start \texttt{ta}; \texttt{ta} is abandoned; start
\texttt{tb} over it (L4 already broken here); \texttt{ta}'s read fails on the
generation both read, which is closed; \texttt{tb} is abandoned; redial; then
\texttt{tb}'s report --- its socket is the dead generation~1 --- closes
generation~2, which was neither closed by the link nor broken by its own
reader.  For L1 it returns the SF3-2 trace: dial, start, feed, the link
closes, and the reader, which no one stopped, routes its generation-1 reply
into the shared queue behind a closed (and, after a redial, replaced)
connection.  For L3 the shortest trace is three steps: start, abandon, the
abandoned reader reports.

\paragraph{The stopping flag as the only repair ---
\texttt{listener\_lifecycle\_gate\_only\_buggy.tex}.}
This is the reviewers' proposed fix taken literally: a reader acts only with
permission, issued while its flag is clear.  The sink is still shared and the
report is still \texttt{end}.  L3 and L4 now hold.  L1 and L2 are
\emph{still violated}: the reader looks, finds the flag clear, is abandoned
and replaced, and acts on the replacement.  The window between the look and
the step is a thread being descheduled, or a handler running, and no
amount of earlier testing of the flag closes it.  Only the confinement of the
step --- a sink and a compare-and-close that name the reader's generation ---
does.

% ============================================================
\section{Verification}
\label{sec:verify}
% ============================================================

Type-checking is a \texttt{fuzz} pass:
\begin{verbatim}
  fuzz docs/listener_lifecycle.tex
\end{verbatim}

Every \texttt{probcli} call below also takes \texttt{-p MAX\_OPERATIONS 200}:
at ProB's default of ten, the operations with an input are truncated in
some states and the check is reported incomplete.  The coverage report
(\texttt{-coverage}) lists every operation as covered; it was the first run
that showed three operations unreachable, which was an error in this file's
frame conditions and has been corrected.

The fixed design.  The negation of each invariant is unreachable (goal
\emph{not found}):
\begin{verbatim}
  # L1 routing stays within a generation
  probcli docs/listener_lifecycle.tex -model_check -nodead \
    -goal "(qgen /= {} & (held = fclear or qgen /= {curGen})) or (waiters /= {} & (held = fclear or waiters /= {curGen}))" \
    -p DEFAULT_SETSIZE 2

  # L2 a generation is closed only by the link or by itself
  probcli docs/listener_lifecycle.tex -model_check -nodead \
    -goal "held = fclear & curGen > 0 & curGen /: (broken \/ deliberate)" \
    -p DEFAULT_SETSIZE 2

  # L3 a stopped reader is silent
  probcli docs/listener_lifecycle.tex -model_check -nodead \
    -goal "leak = fset" -p DEFAULT_SETSIZE 2

  # L4 one reader per live generation
  probcli docs/listener_lifecycle.tex -model_check -nodead \
    -goal "lst(ta) /= lexit & lst(tb) /= lexit & lgen(ta) = lgen(tb)" \
    -p DEFAULT_SETSIZE 2
\end{verbatim}

The scenarios are reachable, so the invariants are not vacuous (goal
\emph{found}): an abandoned reader with its successor running, a reply routed
live, a reply dropped dead, and a stream the reader found genuinely broken.
\begin{verbatim}
  -goal "lst(ta) = lstop & lst(tb) = lrun & lgen(ta) < lgen(tb)"
  -goal "qgen /= {}"
  -goal "broken /= {}"
  -goal "lst(ta) = lstop & buf(ta) > 0 & lgen(ta) < curGen"
\end{verbatim}

Deadlock-freedom and the structural invariants, by the full model-check:
\begin{verbatim}
  probcli docs/listener_lifecycle.tex -model_check -p DEFAULT_SETSIZE 2
\end{verbatim}

The controls make the negations \emph{found}:
\begin{verbatim}
  probcli docs/listener_lifecycle_current_buggy.tex -model_check -nodead \
    -goal "held = fclear & curGen > 0 & curGen /: (broken \/ deliberate)" \
    -p DEFAULT_SETSIZE 2
  probcli docs/listener_lifecycle_gate_only_buggy.tex -model_check -nodead \
    -goal "held = fclear & curGen > 0 & curGen /: (broken \/ deliberate)" \
    -p DEFAULT_SETSIZE 2
\end{verbatim}
(and likewise L1, L3 and L4 against the first, L1 against the second; the
exact outputs are in the coverage note.)

Re-run \texttt{fuzz} and these checks whenever \texttt{listener\_thread.py},
\texttt{display\_replies.py}, \texttt{display\_connection.py} or
\texttt{inbound\_router.py} change how a reader is started, stopped, or
allowed to act on the connection or the reply queues.

\end{document}
